\chapter{Einleitung}

\section{Motivation}
Um Programmieren richtig lernen zu können, muss die Theorie auch in der Praxis angewendet werden. Das wird aber besonders am Anfang durch aufwendiges Herunterladen von Codeeditoren, Einrichten der Programmiersprache und dem System erschwert und schreckt ab.
Dafür wurde an der HTWG Konstanz die Lernplattform CodeTask entwickelt, auf der Studenten strukturiert die Konzepte rund ums Programmieren lernen können. CodeTask setzt dabei auf einfach aufgebaute Aufgaben. Der Großteil der Aufgaben fängt mit einem kurzen Video über die aktuelle Aufgabe an, daraufhin folgt ein kleiner Infotext, in welchem die Aufgabe noch mal erklärt wird. Nach dem Informationstext wird der Student bzw. der Nutzer vor unterschiedliche Aufgabentypen gestellt, die er lösen kann. Darunter auch der CodeTask Aufgabentyp, bei welchem der Student eigenständig in einem Codeeditor auf der Website Code programmieren und validieren lassen kann. Diese Lernplattform sorgt dafür, dass neue Studenten oder Personen, die erste Erfahrungen im Programmieren sammeln, nicht sofort mit überflüssigen Informationen überschüttet werden und sich erst mal auf die wirklich essenziellen Bereiche der Programmierung fokussieren können.
Damit solch ein System funktionieren kann, muss das System in der Lage sein, den Code, welcher von dem Nutzer geschrieben wird, ausführen und gegen vordefinierte Tests prüfen zu können und dem Nutzer ein klares Ergebnis zu liefern, mit dem der Nutzer genau sieht, ob er alles richtig gemacht hat, oder, wenn es ein Problem gab, warum es zu diesem Problem gekommen ist.
CodeTask kann genau das, und zwar durch die Check-API. Ein Service, welcher den Nutzercode zusammen mit den für diese Aufgabe vorgesehen Tests übergeben bekommt, diesen überprüft, ausführt und zum Schluss ein Ergebnis zurück gibt, welches für den Nutzer möglichst verständlich ist.


\section{Problemstellung}
Solch ein System bringt aber einige kritische Probleme mit sich, welche richtig gehandhabt werden müssen, um unerwünschtes Verhalten, wie Systemausfälle oder lange Wartezeiten, zu verhindern. Zu Beginn dieser Arbeit war die Check-API bereits ein eigener Service und behandelte schon einige dieser kritischen Probleme, allerdings nicht alle und manche auch nicht ausreichend genug.
Durch die Entwicklung der Check-API als eigener Service wurde so schon vor dieser Bachelorarbeit das größte Problem sehr stark eingeschränkt. Das Ausführen von Code während der Laufzeit ist für solch ein System eines der größten Sicherheitsrisiken, da willkürlicher Code mit vollständiger Berechtigung auf das System ausgeführt werden kann. Durch die Isolation in einen eigenen Service wird dieses Problem schon stark eingeschränkt, da die Auswirkungen problematischer Codeausführung auf die Check-API begrenzt werden und so die umliegenden Systeme nicht direkt betroffen sind.
Aber nicht nur das Ausführen von Code ist in so einem System ein Problem. Besonders auf den Lernprozess der Nutzer ist es relevant, diese richtig testen zu können. Bisher unterstützte das System lediglich öffentliche Testfälle, welche größtenteils sehr leicht zu umgehen sind und so die eigentliche Lösung nicht wie erwünscht gelöst wird. Des Weiteren bekam der Nutzer nur mitgeteilt, ob die Aufgabe korrekt gelöst wurde oder einen Fehler hatte, dabei wurde aber nicht klar übermittelt, was für ein Fehler es ist und warum er möglicherweise aufgetreten ist.
Als letztes Problem kann die Ausführung von Nutzercode ein Risiko für die Stabilität des Systems darstellen. Programme können ineffizient programmiert oder, schlimmer noch, Endlosschleifen beinhalten. Ohne eine Kontrolle über die Laufzeit kann die Ausführung von solch einem Programm extrem viele Ressourcen verbrauchen oder sogar das gesamte System extrem ausbremsen oder sogar unbrauchbar machen. Problematisch wird so ein Verhalten vor allem, wenn die Plattform von mehreren Nutzern gleichzeitig genutzt wird.

\section{Ziel der Arbeit}
Das Ziel dieser Bachelorarbeit ist die Weiterentwicklung der bestehenden Check-API. Dabei sollen die eben genannten Probleme verbessert oder vollständig gelöst werden. Die Architektur soll dabei nicht ersetzt, sondern gezielt erweitert werden, sodass eine Einreichung von Nutzercode sicher, robust und aussagekräftig verarbeitet werden kann, um eine möglichst effektive und angenehme Lernumgebung zu schaffen.

Ein Schwerpunkt liegt auf der Unterstützung von versteckten Tests. Diese ermöglichen das Überprüfen des Nutzercodes, ohne dass der Nutzer weiß, wie dieser Test aussieht. Dadurch wird erschwert, dass der Nutzer seine Lösung gezielt an die sichtbaren Tests anpasst. Dabei muss beachtet werden, dass der Nutzer zu keinem Zeitpunkt diese versteckten Tests einsehen kann, um auch hier eine gezielte Umgehung zu verhindern.
Des Weiteren ist eine verbesserte Fehlerbehandlung und Ausgabe von Relevanz. Sollte ein Nutzer nur mitgeteilt bekommen, dass ein Fehler vorliegt, aber nicht gesagt wird, warum dieser Fehler aufgetreten ist, verliert der Nutzer die Lust, weil er ohne Weiteres eventuell nicht weiterkommt. Daher ist es wichtig, die möglichen Fehlerquellen zu erkennen und diese auf eine individuell angepasste Weise zu handhaben, um dem Nutzer möglichst klar mitzuteilen, was das Problem ist, ohne dem Nutzer die genaue Lösung vorzulegen.
Darüber hinaus muss die Robustheit weiter verbessert werden. Nicht terminierende oder problematische Einreichungen dürfen den Service nicht permanent außer Betrieb nehmen oder blockieren. Dafür muss vor der Ausführung geprüft werden, ob es in dem eingereichten Code kritische Befehle gibt, welche dem Service schaden können, und zusätzlich muss die Dauer der Ausführung überwacht werden.
Weiterhin wird die Leistung des Systems untersucht, um festzustellen, wie das System schneller und leistungsfähiger gemacht werden kann.

\section{Aufbau der Arbeit}
Kapitel \ref{chap:basicsAndTechnologies} beschreibt die Grundlagen und Technologien, die für das Verständnis der Arbeit notwendig sind. Dazu wird zusätzlich auf vergleichbare Systeme eingegangen, um einen Vergleich mit der Lösung von CodeTask zu schaffen.

In Kapitel \ref{chap:requirements} wird auf das bestehende System sowie die Problemstellung eingegangen. Des Weiteren werden die Anforderungen dieser Arbeit beschrieben. Darunter die funktionalen und nicht funktionalen Anforderungen, die sich für die Check-API und das umliegende System ergeben.

Kapitel \ref{chap:concept&architecture} geht auf das Konzept und die Zielarchitektur der erweiterten Check-API ein. Es geht darauf ein, wie die zuvor definierten Anforderungen umgesetzt werden können und welche Rolle die Check-API in dem System übernimmt.

Die genaue Umsetzung wird in Kapitel \ref{chap:implementierung} erläutert. Dort wird die Anpassung des Parsers zur Unterstützung von versteckten Tests, die Zusammensetzung einer Anfrage an die Check-API, die Umsetzung der Geheimhaltung der versteckten Tests sowie die verbesserte Fehlerbehandlung und Robustheit beschrieben.

Kapitel \ref{chap:evaluation} evaluiert die umgesetzte Lösung. Dafür wird die Fehlerbehandlung aufgeschlüsselt und bewertet, sowie die Check-API in unterschiedlichen Umgebungen auf Performance untersucht und verglichen.

Abschließend fasst Kapitel \ref{chap:fazit} die Ergebnisse dieser Bachelorarbeit zusammen. Zusätzlich werden die Ergebnisse bewertet und ein Ausblick auf mögliche Verbesserungen und Weiterentwicklungen geschaffen.
